\chapter{پاسخِ تمرین‌ها}%
\label{app:solutions}

در این پیوست پاسخِ تمرین‌های منتخبِ هر فصل آمده است. به یاد داشته باشید که در
برنامه‌نویسی اغلب \emph{چند} راهِ درست وجود دارد؛ اگر پاسخِ شما با اینجا فرق
دارد اما درست کار می‌کند، آن هم پاسخی معتبر است. پیش از دیدنِ پاسخ، خودتان
بکوشید.

\section*{فصلِ ۱ شروعِ کار}
\phantomsection
\addcontentsline{toc}{section}{فصلِ ۱}

\textbf{تمرینِ ۱.۱} (نام و سن):
\begin{salamfa}
روال ریشه:
    سرچاپ "نام: سارا"
    سرچاپ "سن: 12"
پایان
\end{salamfa}

\textbf{تمرینِ ۱.۲} (\lr{12 * 8}):
\begin{salamfa}
روال ریشه:
    سرچاپ "12 * 8 =", 12 * 8
پایان
\end{salamfa}

\textbf{تمرینِ ۱.۴} (نام و نام‌خانوادگی در یک خط):
\begin{salamfa}
روال ریشه:
    چاپ "سارا"
    چاپ " "
    سرچاپ "احمدی"
پایان
\end{salamfa}

\section*{فصلِ ۲ متغیرها و گونه‌ها}
\phantomsection
\addcontentsline{toc}{section}{فصلِ ۲}

\textbf{تمرینِ ۲.۱} (شهر، جمعیت، دما):
\begin{salamfa}
روال ریشه:
    شهر : رشته = "اصفهان"
    جمعیت : صحیح۶۴ = 2200000
    دما : اعشار۶۴ = 16.5
    سرچاپ شهر, جمعیت, دما
پایان
\end{salamfa}

\textbf{تمرینِ ۲.۴} (تبدیلِ گونه):
\begin{salamfa}
روال ریشه:
    عدد : صحیح = 9
    اعشاری : اعشار۶۴ = عدد برگردان اعشار۶۴
    سرچاپ عدد, اعشاری
پایان
\end{salamfa}

\section*{فصلِ ۳ عملگرها}
\phantomsection
\addcontentsline{toc}{section}{فصلِ ۳}

\textbf{تمرینِ ۳.۱} (محیط و مساحتِ مستطیل):
\begin{salamfa}
روال ریشه:
    طول : صحیح = 8
    عرض : صحیح = 5
    سرچاپ "مساحت:", طول * عرض
    سرچاپ "محیط:", 2 * (طول + عرض)
پایان
\end{salamfa}

\textbf{تمرینِ ۳.۲} (بخش‌پذیری بر ۳):
\begin{salamfa}
روال ریشه:
    اگر (27 % 3) == 0:
        سرچاپ "بخش‌پذیر است"
    وگرنه:
        سرچاپ "بخش‌پذیر نیست"
    پایان
پایان
\end{salamfa}

\textbf{تمرینِ ۳.۵} (مکعب با \code{\textasciicircum\textasciicircum=}):
\begin{salamfa}
روال ریشه:
    ناپایا ضلع : اعشار۶۴ = 3.0
    ضلع ^^= 3
    سرچاپ ضلع
پایان
\end{salamfa}

خروجی \code{27} است. شروع از \code{ناپایا ضلع : صحیح = 3} خطای کامپایل می‌دهد:
\code{\textasciicircum\textasciicircum} همیشه \code{اعشار۶۴} می‌دهد، و یک \code{اعشار۶۴} بدونِ تبدیلِ
صریح (\code{برگردان صحیح}) در یک متغیرِ \code{صحیح} جا نمی‌گیرد؛ یعنی
\code{ضلع \textasciicircum\textasciicircum= 3} دقیقاً به همان دلیلی رد می‌شود که \code{ضلع = ضلع \textasciicircum\textasciicircum 3} رد
می‌شد.

\section*{فصلِ ۴ شرط‌ها}
\phantomsection
\addcontentsline{toc}{section}{فصلِ ۴}

\textbf{تمرینِ ۴.۱} (مثبت، منفی، صفر):
\begin{salamfa}
روال علامت(ن : صحیح) : رشته:
    اگر ن > 0: برگشت "مثبت" پایان
    اگر ن < 0: برگشت "منفی" پایان
    برگشت "صفر"
پایان

روال ریشه:
    سرچاپ علامت(5), علامت(-3), علامت(0)
پایان
\end{salamfa}

\textbf{تمرینِ ۴.۴} (زوج یا فرد):
\begin{salamfa}
روال زوج است(ن : صحیح) : منطقی:
    برگشت (ن % 2) == 0
پایان

روال ریشه:
    سرچاپ زوج است(4), زوج است(7)
پایان
\end{salamfa}

\section*{فصلِ ۵ حلقه‌ها}
\phantomsection
\addcontentsline{toc}{section}{فصلِ ۵}

\textbf{تمرینِ ۵.۲} (جمعِ زوج‌های ۱ تا ۱۰۰):
\begin{salamfa}
روال ریشه:
    ناپایا مجموع : صحیح = 0
    ناپایا i : صحیح = 2
    تا i <= 100:
        مجموع = مجموع + i
        i = i + 2
    پایان
    سرچاپ مجموع
پایان
\end{salamfa}

\textbf{تمرینِ ۵.۵} (جدولِ ضربِ ۱ تا ۵):
\begin{salamfa}
روال ریشه:
    تکرار 1 تا 5 از i:
        ناپایا خط : رشته = ""
        تکرار 1 تا 5 از j: خط = خط + (i * j) + " " پایان
        سرچاپ خط
    پایان
پایان
\end{salamfa}

\section*{فصلِ ۶ روال‌ها}
\phantomsection
\addcontentsline{toc}{section}{فصلِ ۶}

\textbf{تمرینِ ۶.۳} (فیبوناچی):
\begin{salamfa}
روال فیبوناچی(ن : صحیح) : صحیح:
    اگر ن < 2: برگشت ن پایان
    برگشت فیبوناچی(ن - 1) + فیبوناچی(ن - 2)
پایان

روال ریشه:
    سرچاپ فیبوناچی(10)
پایان
\end{salamfa}

\textbf{تمرینِ ۶.۵} (ب.م.م به روشِ اقلیدس):
\begin{salamfa}
روال بمم(الف : صحیح, ب : صحیح) : صحیح:
    تا ب != 0:
        ناپایا موقت : صحیح = ب
        ب = الف % ب
        الف = موقت
    پایان
    برگشت الف
پایان

روال ریشه:
    سرچاپ بمم(48, 36)
پایان
\end{salamfa}

\section*{فصلِ ۷ آرایه‌ها و بردارها}
\phantomsection
\addcontentsline{toc}{section}{فصلِ ۷}

\textbf{تمرینِ ۷.۱} (بزرگ‌ترینِ آرایه):
\begin{salamfa}
روال ریشه:
    اعداد : صحیح[6] = [12, 45, 7, 23, 56, 9]
    ناپایا بیشینه : صحیح = اعداد[0]
    تکرار 1 تا 5 از i:
        اگر اعداد[i] > بیشینه: بیشینه = اعداد[i] پایان
    پایان
    سرچاپ "بیشینه:", بیشینه
پایان
\end{salamfa}

\textbf{تمرینِ ۷.۳} (زوج‌های یک بردار):
\begin{salamfa}
روال ریشه:
    ناپایا ب : وکتور<صحیح> = وکتور {}
    تکرار 1 تا 10 از i: ب.بیفزا(i) پایان
    ناپایا j : صحیح = 0
    تا j < ب.طول():
        اگر (ب.بگیر(j)[0] % 2) == 0: سرچاپ ب.بگیر(j)[0] پایان
        j = j + 1
    پایان
پایان
\end{salamfa}

\section*{فصلِ ۸ رشته‌ها}
\phantomsection
\addcontentsline{toc}{section}{فصلِ ۸}

\textbf{تمرینِ ۸.۵} (متقارن بودن):
\begin{salamfa}
روال متقارن است(س : رشته) : منطقی:
    ناپایا چپ : صحیح = 0
    ناپایا راست : صحیح = س.طول() - 1
    تا چپ < راست:
        اگر س.زیررشته(چپ, 1) != س.زیررشته(راست, 1):
            برگشت نادرست
        پایان
        چپ = چپ + 1
        راست = راست - 1
    پایان
    برگشت درست
پایان

روال ریشه:
    سرچاپ متقارن است("radar"), متقارن است("salam")
پایان
\end{salamfa}

\section*{فصلِ ۹ ساختارها و شمارشی‌ها}
\phantomsection
\addcontentsline{toc}{section}{فصلِ ۹}

\textbf{تمرینِ ۹.۱} (دایره):
\begin{salamfa}
ساختار دایره:
    شعاع : اعشار۶۴ = 0.0
    روال مساحت() : اعشار۶۴: برگشت 3.1416 * این.شعاع * این.شعاع پایان
    روال محیط() : اعشار۶۴: برگشت 2.0 * 3.1416 * این.شعاع پایان
پایان

روال ریشه:
    ناپایا د : دایره = دایره { شعاع = 5.0 }
    سرچاپ "مساحت:", د.مساحت()
    سرچاپ "محیط:", د.محیط()
پایان
\end{salamfa}

\section*{فصلِ ۱۰ میانجی‌ها}
\phantomsection
\addcontentsline{toc}{section}{فصلِ ۱۰}

\textbf{تمرینِ ۱۰.۲} (جاندار):
\begin{salamfa}
میانجی جاندار:
    روال صدا() : رشته
پایان
ساختار گربه:
    روال صدا() : رشته: برگشت "میو" پایان
پایان
ساختار سگ:
    روال صدا() : رشته: برگشت "واق" پایان
پایان

روال بشنو(ج : dyn جاندار): سرچاپ ج.صدا() پایان

روال ریشه:
    بشنو(گربه {})
    بشنو(سگ {})
پایان
\end{salamfa}

\section*{فصلِ ۱۱ کلی‌سازی}
\phantomsection
\addcontentsline{toc}{section}{فصلِ ۱۱}

\textbf{تمرینِ ۱۱.۱} (کمینهٔ کلی):
\begin{salamfa}
روال کمینه<ت>(الف : ت, ب : ت) : ت:
    اگر الف < ب: برگشت الف پایان
    برگشت ب
پایان

روال ریشه:
    سرچاپ کمینه(3, 9)
    سرچاپ کمینه(2.5, 1.5)
پایان
\end{salamfa}

\section*{فصلِ ۱۲ بستارها و لامبدا}
\phantomsection
\addcontentsline{toc}{section}{فصلِ ۱۲}

\textbf{تمرینِ ۱۲.۲} (سازندهٔ ضرب):
\begin{salamfa}
روال سازنده ضرب(ضریب : صحیح) : func(صحیح) صحیح:
    برگشت (x : صحیح) => x * ضریب
پایان

روال ریشه:
    سه‌برابر : func(صحیح) صحیح = سازنده ضرب(3)
    سرچاپ سه‌برابر(7)
پایان
\end{salamfa}

\section*{فصلِ ۱۳ خطا و دیرکرد}
\phantomsection
\addcontentsline{toc}{section}{فصلِ ۱۳}

\textbf{تمرینِ ۱۳.۲} (ریشهٔ دومِ امن):
\begin{salamfa}
واردسازی "math"
روال ریشه امن(x : اعشار۶۴, نتیجه &: اعشار۶۴) : منطقی:
    اگر x < 0.0: برگشت نادرست پایان
    نتیجه = ریاضی.جذر(x)
    برگشت درست
پایان

روال ریشه:
    ناپایا ر : اعشار۶۴ = 0.0
    اگر ریشه امن(16.0, ر): سرچاپ ر وگرنه: سرچاپ "خطا" پایان
    اگر ریشه امن(-4.0, ر): سرچاپ ر وگرنه: سرچاپ "خطا" پایان
پایان
\end{salamfa}

\section*{فصلِ ۱۴ ماژول‌ها}
\phantomsection
\addcontentsline{toc}{section}{فصلِ ۱۴}

\textbf{تمرینِ ۱۴.۱} (ماژولِ هندسه). فایلِ \code{هندسه.salam}:
\begin{salamfa}
بسته هندسه
همگانی روال مساحت دایره(ش : اعشار۶۴) : اعشار۶۴: برگشت 3.1416 * ش * ش پایان
همگانی روال محیط دایره(ش : اعشار۶۴) : اعشار۶۴: برگشت 2.0 * 3.1416 * ش پایان
\end{salamfa}
فایلِ اصلی:
\begin{salamfa}
واردسازی ه "هندسه.salam"

روال ریشه:
    سرچاپ ه.مساحت دایره(5.0)
    سرچاپ ه.محیط دایره(5.0)
پایان
\end{salamfa}

\section*{فصلِ ۱۵ کتابخانهٔ استاندارد}
\phantomsection
\addcontentsline{toc}{section}{فصلِ ۱۵}

\textbf{تمرینِ ۱۵.۱} (فیثاغورس):
\begin{salamfa}
واردسازی "math"

روال ریشه:
    a : اعشار۶۴ = 3.0
    b : اعشار۶۴ = 4.0
    سرچاپ "وتر:", ریاضی.جذر((a * a) + (b * b))
پایان
\end{salamfa}

\textbf{تمرینِ ۱۵.۲} (نگاشتِ قیمت‌ها):
\begin{salamfa}
روال ریشه:
    ناپایا قیمت : نگاشت<رشته, صحیح> = نگاشت {}
    قیمت.درج("سیب", 12000)
    قیمت.درج("موز", 18000)
    سرچاپ "قیمتِ سیب:", قیمت.بگیر("سیب")
پایان
\end{salamfa}

\section*{فصلِ ۱۶ فایل و سیستم}
\phantomsection
\addcontentsline{toc}{section}{فصلِ ۱۶}

\textbf{تمرینِ ۱۶.۱} (نوشتن و خواندنِ فایل):
\begin{salamen}
package main
import "io"

func main:
    io.WriteFile("note.txt", "سلام از فایل!")
    text : str = io.ReadFile("note.txt")
    println text
end
\end{salamen}

\section*{فصلِ ۱۷ هم‌روندی}
\phantomsection
\addcontentsline{toc}{section}{فصلِ ۱۷}

\textbf{تمرینِ ۱۷.۳} (چرا \code{local}؟): اگر هر نخ در هر گام مستقیم به
\code{sum} اضافه می‌کرد، یا باید در هر گام قفل می‌گرفت (که هم‌روندی را بی‌اثر
می‌کرد)، یا بدونِ قفل دچارِ \emph{رقابت} می‌شد و نتیجه نادرست در می‌آمد. با
جمعِ محلی، هر نخ آزادانه و بدونِ قفل کار می‌کند و تنها \emph{یک‌بار} در پایان
قفل می‌گیرد هم سریع و هم درست.

\section*{فصلِ ۱۸ ارتباط با کدِ بیرونی}
\phantomsection
\addcontentsline{toc}{section}{فصلِ ۱۸}

\textbf{تمرینِ ۱۸.۱} (\code{sqrt} از C):
\begin{salamen}
extern:
    func printf(format : str, ...) : int
    func sqrt(x : f64) : f64
end

func main:
    printf("sqrt(2)  = %.4f\n", sqrt(2.0))
    printf("sqrt(16) = %.4f\n", sqrt(16.0))
end
\end{salamen}

\section*{فصلِ ۱۹ وب و چیدمان}
\phantomsection
\addcontentsline{toc}{section}{فصلِ ۱۹}

\textbf{تمرینِ ۱۹.۲} (فهرستِ کارها):
\begin{salamfa}
چیدمان:
    عنوان = "کارهای امروز"
    جهت = "راست به چپ"
    فهرست:
        مورد: محتوا = "خواندنِ کتابِ سلام" پایان
        مورد: محتوا = "نوشتنِ یک برنامه" پایان
        مورد: محتوا = "حلِ تمرین‌ها" پایان
    پایان
پایان
\end{salamfa}

\section*{فصلِ ۲۰ نکاتِ پیشرفته}
\phantomsection
\addcontentsline{toc}{section}{فصلِ ۲۰}

\textbf{تمرینِ ۲۰.۱} (جابجاییِ دو رشته با ارجاع):
\begin{salamfa}
روال جابجایی(الف &: رشته, ب &: رشته):
    ناپایا موقت : رشته = الف
    الف = ب
    ب = موقت
پایان

روال ریشه:
    ناپایا x : رشته = "اول"
    ناپایا y : رشته = "دوم"
    جابجایی(x, y)
    سرچاپ x, y
پایان
\end{salamfa}

\vspace{1cm}
\begin{center}
\emph{به اینجا که رسیدید، دیگر یک برنامه‌نویسِ سلام هستید. راهِ یادگیری هرگز
پایان نمی‌یابد پس بنویسید، بیازمایید و لذت ببرید!}
\end{center}
